% Publication figure snippets for FLOW-CM

\begin{figure}[H]
\centering
\includegraphics[width=0.92\textwidth]{figures/model_pipeline.pdf}
\caption{Schematic structure of FLOW-CM. The baryonic distribution generates nodal occupation, structural collectivity, inward memory, basal connectivity, radial transition, graph current, harmonic conductance, and finally the predicted rotation curve.}
\label{fig:pipeline}
\end{figure}

\begin{figure}[H]
\centering
\includegraphics[width=0.98\textwidth]{figures/rotation_curve_examples.pdf}
\caption{Representative rotation curves comparing observed velocities, baryonic Newtonian predictions, the RAR/MOND benchmark, and harmonic FLOW-CM.}
\label{fig:curves}
\end{figure}

\begin{figure}[H]
\centering
\includegraphics[width=0.78\textwidth]{figures/mse_ratio_histogram.pdf}
\caption{Distribution of per-galaxy MSE ratios between FLOW-CM and the external RAR/MOND benchmark. Values below unity indicate galaxies for which FLOW-CM gives a lower velocity MSE.}
\label{fig:mse_ratio}
\end{figure}

\begin{figure}[H]
\centering
\includegraphics[width=0.72\textwidth]{figures/permutation_test.pdf}
\caption{Geometric permutation test. The increase in MSE after permuting $Q_{\rm harm}$ between galaxies indicates that the radial nodal conductance contains galaxy-specific information.}
\label{fig:perm}
\end{figure}
